\scanpage{189}

\[\xi^n=L_n(E_1,\ldots,E_n)=\sum_{i=1}^{n}c_{ni}E_i\qquad\text{\uuline{$c_{nn}=1$}}\]

On en conclut la

\noindent\begin{tabular}{@{}|p{.92\linewidth}@{}}\uline{Proposition 3} $\unsure{A|\mathcal A}(k)$ est l'anneau de polynômes $Z[\xi]$ sur $Z$ engendré par $\xi=E_2$ $(d(\xi)=2)$\end{tabular}

N. B. Le calcul explicite donne

\diagram{\(\begin{array}{l@{\qquad\qquad}l}\xi^0=E_1&E_1=1\\\xi^1=E_2&E_2=\xi\\\xi^2=E_1+E_3&E_3=\xi^2-1\\\xi^3=2E_2+E_4&E_4=\xi^3-2\xi\\\xi^4=2E_1+3E_3+E_5&E_5=\xi^4-3\xi^2+5\\\ldots&\ldots\end{array}\)}

(Toutes les structures de $A(k)$ sont bien précisées dans \unsure{ce|un} th. de structure, \uline{sauf} la structure d'ordre, qui est un peu cachée !)

% archive sheet number: 188
